%! Absolute Value Functions \& Transformations — Unit 1, Lesson 1.3 — Student Packet Cover Sheet
% Regenerated 2026-09-07 in the course-wide gradual-release shape (LESSON_SHAPE.md §1, §2):
% three scored rows, formal vocabulary in the targets, a content summary in Keep in Mind, and the
% course's due-date policy on the homework row.
\documentclass[10pt]{article}
\usepackage{algebra2-article}
\usepackage{algebra2-boxes}
\usepackage{ltablex}
\keepXColumns

\begin{document}

% Full-bleed forest banner
% The header text is set INSIDE a tikz node anchored to the page so it can never
% drift above the paper edge (a negative \vspace here pushes it off the sheet).
\begin{tikzpicture}[remember picture, overlay]
  \fill[forest] (current page.north west)
    rectangle ([yshift=-1.16in]current page.north east);
  \node[anchor=north, align=center, inner sep=0pt, text width=\paperwidth]
    at ([yshift=-0.28in]current page.north) {%
      {\color{white}\bfseries\LARGE Algebra 2}\\[4pt]
      {\color{forestmid}\bfseries\large Unit 1 \quad Linear Functions}\\[6pt]
      {\color{white}\bfseries\large Lesson 1.3 \quad Absolute Value Functions \& Transformations}%
    };
\end{tikzpicture}

% Push the body clear of the 1.16in banner (page top margin is 0.60in).
\vspace*{0.56in}
\namedateperiod
\vspace{0.12in}

% Learning targets use the lesson's FORMAL VOCABULARY — the Guided Notes name every term before
% the practice uses it, so there is nothing to withhold.
\begin{learningtargetbox}
\small
\begin{enumerate}[leftmargin=1.5em, itemsep=5pt]
  \item \ldots{} read the \textbf{absolute value parent} $f(x)=|x|$ as a V, and name its
        \textbf{vertex}, \textbf{axis of symmetry}, and \textbf{minimum}.
  \item \ldots{} use \textbf{vertex form} $g(x)=a|x-h|+k$ to find the vertex $(h,k)$ and the axis
        $x=h$ --- and explain why the inside of the bars moves the graph \emph{backwards}.
  \item \ldots{} decide from $a$ whether a V opens up (a \textbf{minimum}) or down (a
        \textbf{reflection}, so $k$ is a \textbf{maximum}), and whether it is \textbf{narrower} or
        \textbf{wider} than the parent.
  \item \ldots{} read domain, range, and increasing/decreasing intervals off a V, build its rule
        from a graph, and connect the two solutions of $|ax+b|=c$ to where the V crosses $y=c$.
\end{enumerate}
\end{learningtargetbox}

\vspace{0.14in}

\begin{tocbox}
\small
\renewcommand{\arraystretch}{1.5}
\begin{tabularx}{\linewidth}{@{} c l X r @{}}
\rowcolor{forestbg}
\textbf{\#} & \textbf{Component} & \textbf{Description} & \textbf{Score} \\
\midrule
% Rows are in packet order (shared/lesson.mk STUDENT_ORDER). THREE scored rows: no group
% activity, no exit ticket, and the debrief has no row (a phase, not a component). Every row is
% scored — a \blank{1.2cm}, never NA. Four columns (c l X r): every row needs FOUR cells.
1 & Warm-Up        & Spiral review: a table for $y=|x|$, the three moves from Lesson 1.1, and
                     reading a V                                                                 & \blank{1.2cm} \\
2 & Guided Notes   & The parent V; inside vs.\ outside the bars; the sign and size of $a$;
                     vertex form all at once --- ending in \emph{Guided Practice} and
                     \emph{Individual Practice}                                                  & \blank{1.2cm} \\
3 & Homework       & Features from three rules, a one-sign-apart pair, matching graphs and a
                     table, the odd cases, a drone model, and one multiple-choice item ---
                     \textbf{scored; due the first class after two study halls}                  & \blank{1.2cm} \\
\midrule
  & \multicolumn{2}{r}{\textbf{Total}} & \blank{1.2cm} \\
\end{tabularx}
\end{tocbox}

\vspace{0.10in}

% Keep in Mind is a CONTENT summary — the definitions and the distinction the lesson turns on.
\begin{remindbox}
\small The graph of $y=|x|$ is a \textbf{V} with its \textbf{vertex} at $(0,0)$, because a distance
can never dip below zero --- so the vertex is a turning point and the graph is a mirror image
across its \textbf{axis of symmetry}. In \textbf{vertex form} $g(x)=a|x-h|+k$ the vertex is
$(h,k)$: find $h$ by setting the inside of the bars to zero (\emph{inside is horizontal and
backwards}), and $k$ moves the graph up or down. The sign of $a$ says \emph{up} ($k$ is a
\textbf{minimum}) or \emph{down} ($k$ is a \textbf{maximum} --- a negative $a$ flips the vertex's
\emph{job}, not just the picture); the size of $a$, the slope of the right-hand ray, says
\emph{narrower} ($|a|>1$) or \emph{wider} ($0<|a|<1$).
\end{remindbox}

\end{document}
